\documentclass[11pt]{article}
\usepackage[margin=1in]{geometry}
\usepackage{graphicx}
\usepackage{booktabs}
\usepackage{amsmath}
\usepackage{amssymb}
\usepackage{siunitx}
\usepackage{hyperref}
\usepackage{natbib}
\usepackage{multirow}

\title{MANTA: Physics-Derived Neural Architecture for Exoplanet Transit Detection}
\author{Author One \and Author Two \and Author Three}
\date{\today}

\begin{document}
\maketitle

\begin{abstract}
We present MANTA (Mandel-Agol Neural Transit Architecture), a transit detection
network whose architecture is derived directly from transit physics rather than
empirical neural design search. MANTA encodes five constraints from the
Mandel-Agol formalism and stellar variability physics: position-aware attention
from limb darkening, smooth non-monotonic activation from elliptic-integral
ingress/egress geometry, bounded output from physical flux constraints,
palindromic convolution from ingress/egress time reversibility, and parallel
frequency streams for granulation/asteroseismology/starspot separation. On
Kepler DR25, MANTA achieves competitive or improved performance versus AstroNet
with improved physical interpretability and calibration.
\end{abstract}

\section{Introduction}
\subsection{Motivation}
Exoplanet transit classification is a high-volume astrophysical inference
problem where model inductive bias can materially affect detectability of small
and grazing planets.

\subsection{Contribution}
MANTA introduces architecture-as-physics: constraints are encoded directly in
network structure, not only in training objectives.

\section{Background and Related Work}
\subsection{Transit Physics}
The Mandel-Agol model describes transit flux under stellar limb darkening and
planetary occultation geometry \citep{mandelagol2002}.

\subsection{Neural Transit Detection}
AstroNet established deep-learning transit classification at scale
\citep{shalluevanderburg2018}. ExoMiner advanced production vetting
\citep{valizadegan2021}. Recent work includes GPFC and ExoSpikeNet
\citep{kunimoto2024gpfc, exospikenet2024}.

\section{Physics-Derived Constraints}
\subsection{Constraint 1: Limb Darkening to Position-Aware Attention}
MANTA's attention bias is not merely position-dependent in sequence index.
Instead, each phase position in the folded transit window is mapped to a
projected radius on the stellar disk along a transit chord parameterized by
impact parameter $b$. Let $s \in [-1,1]$ denote normalized phase position along
the transit window. We map $s$ to projected radius
\begin{equation}
  r(s,b) = \sqrt{b^2 + (1-b^2)s^2},
\end{equation}
then evaluate the quadratic limb-darkening law
\begin{equation}
  I(r) = 1 - u_1\left(1-\mu\right) - u_2\left(1-\mu\right)^2,
  \quad \mu = \sqrt{1-r^2}.
\end{equation}
The resulting intensity profile is converted into an additive attention-logit
bias, so attention routing is explicitly chord-geometry-dependent. This ties
the network's inductive bias to the physically relevant distinction between
central ($b \approx 0$) and grazing ($b \rightarrow 1$) transit chords.
\subsection{Constraint 2: Elliptic Integrals to Smooth Activation}
\subsection{Constraint 3: Flux Bounds to Physics-Constrained Output}
\subsection{Constraint 4: Time Reversal to Palindromic Convolution}
\subsection{Constraint 5: Stellar Bands to Parallel Frequency Streams}
Constraint 5 is implemented on the raw temporal light curve, where stellar
variability frequencies are physically meaningful. We first decompose the
cleaned raw time series into granulation, asteroseismology, and starspot bands
in the frequency domain. Each decomposed band is then phase-folded
independently using the target ephemeris, and these folded band-specific views
are passed to their corresponding neural branches.

This ordering is essential: applying FFT directly to an already phase-folded
profile is physically nonsensical for stellar-variability separation, because
phase folding destroys the original temporal frequency structure and aliases it
into morphological (folded-shape) frequencies. Therefore, all frequency-band
interpretation in MANTA is defined pre-fold in the raw temporal domain.

\section{MANTA Architecture}
\subsection{System Overview}
\begin{figure}[t]
  \centering
  \fbox{\parbox{0.9\linewidth}{Placeholder: Full MANTA architecture diagram.}}
  \caption{MANTA architecture with five physics-derived components.}
  \label{fig:architecture}
\end{figure}

\subsection{Implementation Details}
Include layer dimensions, parameter counts, and training/inference complexity.

\section{Experimental Setup}
\subsection{Dataset}
Kepler DR25 TCE catalog with star-level stratified splits and class-imbalance
handling.

\subsection{Training Protocol}
Focal BCE, warmup-cosine scheduling, early stopping, and multi-seed reporting.

\section{Results}
\subsection{Main Benchmark}
\begin{table}[t]
  \centering
  \begin{tabular}{lcccccc}
    \toprule
    Model & Params & AUC-ROC & AUC-PR & F1 & AP & ECE \\
    \midrule
    AstroNet & -- & -- & -- & -- & -- & -- \\
    MANTA & -- & -- & -- & -- & -- & -- \\
    \bottomrule
  \end{tabular}
  \caption{Primary benchmark comparison on Kepler DR25 test split.}
  \label{tab:main_results}
\end{table}

\subsection{Ablation Study}
\begin{table}[t]
  \centering
  \begin{tabular}{lccc}
    \toprule
    Variant & AUC-ROC & F1 & AP \\
    \midrule
    Full MANTA & -- & -- & -- \\
    Minus Position Attention & -- & -- & -- \\
    Minus Elliptic Activation & -- & -- & -- \\
    Minus Symmetric Encoder & -- & -- & -- \\
    Minus Frequency Separation & -- & -- & -- \\
    Minus Physics Output & -- & -- & -- \\
    AstroNet Baseline & -- & -- & -- \\
    \bottomrule
  \end{tabular}
  \caption{Ablation table template for component-wise contribution analysis.}
  \label{tab:ablation}
\end{table}

\subsection{Figures}
\begin{figure}[t]
  \centering
  \fbox{\parbox{0.9\linewidth}{Placeholder: Figure 1 transit detection example.}}
  \caption{Example transit detection and model confidence.}
  \label{fig:transit_detection}
\end{figure}

\begin{figure}[t]
  \centering
  \fbox{\parbox{0.9\linewidth}{Placeholder: Figure 2 frequency decomposition.}}
  \caption{Physics-motivated decomposition into granulation, asteroseismology,
and starspot bands.}
  \label{fig:frequency_decomposition}
\end{figure}

\begin{figure}[t]
  \centering
  \fbox{\parbox{0.9\linewidth}{Placeholder: Figure 3 activation comparison.}}
  \caption{EllipticMish versus Mish and ReLU.}
  \label{fig:activation_comparison}
\end{figure}

\begin{figure}[t]
  \centering
  \fbox{\parbox{0.9\linewidth}{Placeholder: Figure 4 attention over stellar disk.}}
  \caption{Position-aware attention aligned with limb-darkened geometry.}
  \label{fig:attention_weights}
\end{figure}

\begin{figure}[t]
  \centering
  \fbox{\parbox{0.9\linewidth}{Placeholder: Figure 5 ROC curves.}}
  \caption{ROC curves for MANTA, AstroNet, and ablation variants.}
  \label{fig:roc_curves}
\end{figure}

\begin{figure}[t]
  \centering
  \fbox{\parbox{0.9\linewidth}{Placeholder: Figure 6 ablation heatmap.}}
  \caption{Ablation-performance heatmap.}
  \label{fig:ablation_heatmap}
\end{figure}

\begin{figure}[t]
  \centering
  \fbox{\parbox{0.9\linewidth}{Placeholder: Figure 7 planet-size performance.}}
  \caption{Per-planet-size performance profile.}
  \label{fig:planet_size}
\end{figure}

\begin{figure}[t]
  \centering
  \fbox{\parbox{0.9\linewidth}{Placeholder: Figure 8 calibration curves.}}
  \caption{Model calibration and reliability analysis.}
  \label{fig:calibration}
\end{figure}

\section{Discussion}
Discuss where physics-derived constraints provide largest gains and where
limitations remain (e.g., cadence, stellar class coverage).

\section{Conclusion}
Summarize architecture-as-physics framing and implications for scientific ML.

\bibliographystyle{apalike}
\bibliography{references}

\appendix
\section{Appendix A: Additional Derivations}
\section{Appendix B: Hyperparameter Sensitivity}
\section{Appendix C: Additional Tables/Figures}

\end{document}
